\section{Related Work}
\label{sec:related}

Our work sits at the intersection of Symbolic Regression, Simulation-Based Inference, and LLM-driven code generation, building upon decades of research in automated science.

\subsection{Classical and Neural Symbolic Regression}
Symbolic regression seeks to find mathematical expressions that best fit a given dataset while remaining as simple and interpretable as possible. Classical methods rely heavily on Genetic Programming (GP). Algorithms like Eureqa and modern implementations like PySR represent equations as expression trees and evolve them through crossover and mutation. While highly successful in physics environments with clean data, GP suffers from a massive combinatorial explosion when scaling to high-dimensional biological or macroeconomic systems. The search space is dominated by meaningless combinations (e.g., adding a temperature to a logarithmic currency value), and evaluating every candidate is computationally prohibitive.

An alternative approach is Sparse Identification of Non-linear Dynamics (SINDy), which assumes a predefined library of basis functions (e.g., $\sin(x), x^2, e^x$) and uses sparse regression (like LASSO) to select the active terms. SINDy is computationally efficient but fails catastrophically when the true physical mechanism is not in the library or requires a novel boundary condition, such as a localized sigmoidal shift or a stochastic jump.

\subsection{Neural Differential Equations and Black-Box Modeling}
To avoid the combinatorial constraints of symbolic regression, Neural Ordinary Differential Equations (Neural ODEs) and Physics-Informed Neural Networks (PINNs) parameterize the derivative function using a deep neural network. While these methods achieve state-of-the-art interpolation and can be embedded within physical constraints (as in PINNs), they remain fundamentally black boxes. A Neural ODE cannot output a readable formula; it cannot confirm whether a biological process is driven by Michaelis-Menten kinetics or simple mass-action. They are highly susceptible to overfitting noisy data and struggle with stiff systems where local gradients vanish.

\subsection{Simulation-Based Inference (SBI) and ABC-SMC}
Simulation-Based Inference (SBI) bypasses the need for explicit likelihood functions by comparing simulated data directly to observations. This is critical for complex dynamical systems where the likelihood is intractable. Approximate Bayesian Computation Sequential Monte Carlo (ABC-SMC) is a classical SBI algorithm that approximates the posterior distribution of parameters by simulating candidate systems, applying a distance metric, and sequentially filtering particles through decreasing tolerance thresholds.

While traditional SBI focuses solely on parameter estimation for a \textit{static} model structure, $E^3$ innovates by using ABC-SMC to evaluate \textit{dynamically generated} model architectures. The engine treats the LLM as a structural mutator and the ABC-SMC as an unforgiving, statistically rigorous fitness function, seamlessly updating the LLM's prompt context with posterior parameter metrics.

\subsection{Large Language Models in Code Generation and Discovery}
Recent advances have utilized LLMs to generate hypotheses, write Python scripts, or assist in theorem proving. Systems like AlphaCode demonstrate unprecedented capability in syntactic code generation for competitive programming. More recently, DeepMind's FunSearch \citep{romera2023mathematical} and AlphaGeometry \citep{trinh2024solving} have shown that LLMs, when coupled with an automated evaluator, can discover new mathematical bounds and prove complex theorems. 

However, these breakthroughs have largely been confined to discrete mathematics, combinatorics, and formal geometry. In the realm of physical sciences, LLMs typically suffer from ``hallucination,'' proposing plausible-looking but physically incorrect or numerically unstable equations. $E^3$ studies this gap by coupling the LLM with a continuous dynamical simulator (JAX/Diffrax) and an empirical ABC-SMC evaluator. This forces the model to iteratively refine its continuous dynamical hypotheses against empirical data.
